\subsection*{AI 使用声明}

Codex 中的 GPT-5.5 协助生成了两个评测轨道的候选合成场景。研究者在纳入数据集前对候选样本进行了审核和修改，具体过程见第~\ref{sec:construction} 节。大语言模型还按照第~\ref{sec:scoring} 节和附录~\ref{sec:supp-evaluation} 所述协议担任自动裁判。生成式人工智能工具协助将中文稿翻译成英文、起草和修改段落、检索文献和核查引用、完善研究问题与结果解释，以及修改代码和图表。作者负责核实 AI 辅助产生的材料，并对最终论文、数据集、代码、分析和结论的准确性与完整性承担责任。

\subsection*{伦理声明}

本基准使用合成场景，不采集真实青少年的对话。研究者审核了场景与预设风险的对应关系、情境合理性和语言质量。数据包含敏感情境及可能有害的模型回答，用于研究安全失误。这些材料面向安全评测与研究，不应作为建议直接提供给青少年，也不应用于实施有害行为。人工评分抽查由一名标注员完成，具体过程见附录~\ref{sec:supp-human-audit}。基准分数描述模型对所构建场景的回答表现，不能据此认定模型适合在缺乏监督的情况下供青少年使用。

\subsection*{可复现性声明}

第~\ref{sec:construction} 节介绍场景构建与审核流程。第~\ref{sec:scoring} 节和附录~\ref{sec:supp-evaluation} 说明评分协议、裁判输入、人工抽查和统计方法。附录~\ref{sec:supp-results} 给出模型标识及完整结果。附录~\ref{sec:supp-release} 说明复现实验所需的记录，以及通过第三方端点识别自动裁判模型的局限。
